\documentclass[11pt,a4paper]{article}

\usepackage[utf8]{inputenc}
\usepackage[T1]{fontenc}
\usepackage[spanish,es-nodecimaldot]{babel}
\usepackage{lmodern}
\usepackage{microtype}
\usepackage{xcolor}
\usepackage{amssymb}
\usepackage[a4paper,margin=2.4cm]{geometry}
\usepackage[hidelinks]{hyperref}
\usepackage{enumitem}
\usepackage{fancyhdr}

\definecolor{rojocambio}{HTML}{B00020}
\definecolor{grislinea}{HTML}{555555}
\newcounter{cambio}

\newcommand{\camb}[1]{\textcolor{rojocambio}{\textit{#1}}}
\newcommand{\propuesta}[6]{%
  \refstepcounter{cambio}%
  \subsection*{Cambio \thecambio\enspace---\enspace #1}%
  \noindent{\color{grislinea}\texttt{\detokenize{#2}}, #3.}\par
  \smallskip
  \noindent\textbf{Motivo.} #4\par
  \smallskip
  \noindent\textbf{Original.}\par
  \begin{quote}\normalfont #5\end{quote}
  \noindent\textbf{Sugerencia.}\par
  \begin{quote}\normalfont #6\end{quote}
  \noindent\textcolor{grislinea}{$\square$ Aprobar\qquad $\square$ Rechazar\qquad $\square$ Revisar}\par
  \medskip\hrule\medskip
}

\pagestyle{fancy}
\fancyhf{}
\fancyhead[L]{\textit{Leer la Odisea en tiempos iletrados}}
\fancyhead[R]{Informe editorial}
\fancyfoot[C]{\thepage}
\setlength{\headheight}{14pt}
\setcounter{secnumdepth}{0}
\setcounter{tocdepth}{1}
\setlength{\parindent}{0pt}
\setlength{\parskip}{0.5em}

\title{\textbf{Informe editorial de\\
\textit{Leer la Odisea en tiempos iletrados}}}
\author{}
\date{23 de agosto de 2026}

\begin{document}
\maketitle

\section{Criterio de la revisión}

Este informe revisa los veinte capítulos activos que carga \texttt{odisea\_book.tex}. Las líneas indicadas son las líneas físicas de los archivos fuente en el estado del 23 de agosto de 2026. No se ha modificado ninguno de esos archivos.

Se han buscado erratas, problemas ortográficos o sintácticos, repeticiones léxicas próximas y reiteraciones de una idea ya formulada. Se excluyen del escrutinio los pasajes comentados en LaTeX. Las citas de Homero se respetan como texto de Rodríguez Tobal y no se normalizan, aunque contengan giros deliberadamente arcaizantes. Tampoco se consideran defectos las anáforas, los estribillos, las enumeraciones enfáticas ni las fórmulas homéricas cuando cumplen una función rítmica o estructural.

En cada propuesta, el original aparece en redonda. La sugerencia reproduce el pasaje y marca en \camb{cursiva roja} exclusivamente lo que cambia. Las observaciones rotuladas como \emph{repetición} son más discutibles que las erratas y quedan expresamente sometidas al criterio del autor.

\tableofcontents
\clearpage

\section{Introducción: De libros y cine}

\propuesta{Concordancia verbal}{src/odisea0_intro.tex}{línea 71}{La construcción condicional en imperfecto pide imperfecto de subjuntivo; \emph{se perpetúe} rompe la correlación temporal.}{Claro que también sería deseable que la voluntad pedagógica que mostraron Rossi y la RAI se perpetúe, sobre todo si tenemos en cuenta los argumentos sobre los que venimos reflexionando desde hace un rato.}{Claro que también sería deseable que la voluntad pedagógica que mostraron Rossi y la RAI \camb{se perpetuara}, sobre todo si tenemos en cuenta los argumentos sobre los que venimos reflexionando desde hace un rato.}

\propuesta{Repetición léxica}{src/odisea0_intro.tex}{línea 79}{\emph{Traducción} aparece tres veces en el párrafo. La segunda puede sustituirse por un pronombre sin perder precisión.}{La película se estrenó en julio de 2026, a la vez que salía a imprenta la maravillosa y atrevida traducción de la \emph{Odisea} de Juan Manuel Rodríguez Tobal, publicada por Hiperión. Confieso que la lectura de esa traducción fue, literalmente, un empujón al túnel del tiempo.}{La película se estrenó en julio de 2026, a la vez que salía a imprenta la maravillosa y atrevida traducción de la \emph{Odisea} de Juan Manuel Rodríguez Tobal, publicada por Hiperión. Confieso que \camb{su lectura} fue, literalmente, un empujón al túnel del tiempo.}

\section{Primera parte: El cine teme a Homero}

\propuesta{Puntuación}{src/odisea1_nolan.tex}{línea 127}{La conjunción \emph{ni} coordina dos complementos de \emph{hablarnos}; la coma interpuesta no es necesaria.}{Ignora cómo Polifemo ---nada menos que el hijo de un dios--- se pasa la \emph{xenía} por el forro y no se molesta en hablarnos de un personaje central, Nausícaa, ni de su gente, los feacios, a pesar de que estos se caracterizan, precisamente, por respetarla a un elevado coste.}{Ignora cómo Polifemo ---nada menos que el hijo de un dios--- se pasa la \emph{xenía} por el forro y no se molesta en hablarnos de un personaje central, Nausícaa\camb{ ni} de su gente, los feacios, a pesar de que estos se caracterizan, precisamente, por respetarla a un elevado coste.}

\section{Segunda parte: Una osada Odisea}

\propuesta{Ortografía de nombre propio}{src/odisea2_osada.tex}{línea 5}{El programa ferroviario europeo se escribe \emph{Interrail}, con mayúscula inicial.}{Luego hicimos algunos otros recados, preparando su inminente viaje en tren por Europa, el interrail bendito que une los veranos de tantos jóvenes del continente.}{Luego hicimos algunos otros recados, preparando su inminente viaje en tren por Europa, el \camb{Interrail} bendito que une los veranos de tantos jóvenes del continente.}

\propuesta{Puntuación de inciso}{src/odisea2_osada.tex}{línea 13}{La locución incidental \emph{cómo no} debe quedar delimitada por comas.}{Y cómo no, el libro me lo quedé yo (en préstamo).}{Y\camb{,} cómo no, el libro me lo quedé yo (en préstamo).}

\propuesta{Mayúscula dentro de una misma oración}{src/odisea2_osada.tex}{líneas 227--233}{La estrofa nueva no inicia una oración nueva: el verso anterior termina en coma y la subordinada continúa.}{Ya nada más objeto, compañero,\\
Vayacondios, por el ponto mi Ulises,\\
que es vagabundo, mañoso y artero,

Que es de todos los hombres el primero\\
y el último también, hombre de ardides,\\
como todo hombre, falso y verdadero.}{Ya nada más objeto, compañero,\\
Vayacondios, por el ponto mi Ulises,\\
que es vagabundo, mañoso y artero,

\camb{que} es de todos los hombres el primero\\
y el último también, hombre de ardides,\\
como todo hombre, falso y verdadero.}

\section{Tercera parte: Helena}

\propuesta{Ortografía de título}{src/odisea3_helena.tex}{línea 73}{En español solo lleva mayúscula la primera palabra del título. El resto del libro emplea ya la forma correcta.}{El astuto director desea evitar que la película se le deslice del lado de \emph{Harry Potter} o \emph{El Señor de los Anillos} y prefiere salir del paso con cuatro menciones indirectas y una proyección psicológica del síndrome postraumático del guerrero.}{El astuto director desea evitar que la película se le deslice del lado de \emph{Harry Potter} o \emph{El \camb{señor de los anillos}} y prefiere salir del paso con cuatro menciones indirectas y una proyección psicológica del síndrome postraumático del guerrero.}

\section{Menelao}

\propuesta{Puntuación de inciso}{src/odisea3_menelao_es.tex}{líneas 155--158}{La locución \emph{de hecho} es incidental y pide coma.}{Y así, Menelao no habla de Helena,\\
hace ya tantos años que se fue,\\
de hecho desapareció poco después\\
de la visita de Telémaco a Esparta.}{Y así, Menelao no habla de Helena,\\
hace ya tantos años que se fue,\\
de hecho\camb{,} desapareció poco después\\
de la visita de Telémaco a Esparta.}

\propuesta{Cierre de inciso}{src/odisea3_menelao_es.tex}{línea 200}{El inciso que compara dos motivaciones debe cerrarse antes de coordinar el segundo verbo principal.}{Odiseo reflexiona sobre el primer rescate de Helena, más motivado quizás por la codicia del botín que por la necesidad de recuperar a la adúltera y muestra la futilidad de toda la empresa.}{Odiseo reflexiona sobre el primer rescate de Helena, más motivado quizás por la codicia del botín que por la necesidad de recuperar a la adúltera\camb{,} y muestra la futilidad de toda la empresa.}

\propuesta{Construcción correlativa}{src/odisea3_menelao_es.tex}{línea 253}{La serie \emph{ni..., ni..., sino...} queda sin primer término negativo y choca con \emph{al más humilde} colocado antes de ella.}{La mía aquí le hace elegir al más humilde de los héroes, ni al hermoso ---y cobarde--- Paris, ni al aguerrido ---e inútilmente enamorado--- Menelao, sino al callado, gentil Deífobo.}{La mía aquí le hace elegir \camb{no al hermoso ---y cobarde--- Paris, ni al aguerrido ---e inútilmente enamorado--- Menelao, sino al más humilde de los héroes, el callado, gentil Deífobo}.}

\section{Cuarta parte: Calipso}

\propuesta{Repetición léxica}{src/odisea4_calipso.tex}{línea 44}{\emph{Falta} se repite dos veces dentro de una enumeración muy breve.}{La claridad con que Homero desvela su hipocresía, su falta de escrúpulos y su falta de compasión es apabullante.}{La claridad con que Homero desvela su hipocresía, su falta de escrúpulos y su \camb{ausencia} de compasión es apabullante.}

\propuesta{Repetición léxica}{src/odisea4_calipso.tex}{línea 48}{\emph{Puesta en escena} y \emph{escenas} quedan demasiado próximas.}{En el caso del cine, el problema es todavía peor, porque cualquier puesta en escena de las deidades evoca en los espectadores las escenas de acción de mil películas de espada y brujería, o las versiones Marvel de esos mismos dioses.}{En el caso del cine, el problema es todavía peor, porque cualquier \camb{representación} de las deidades evoca en los espectadores las escenas de acción de mil películas de espada y brujería, o las versiones Marvel de esos mismos dioses.}

\section{Quinta parte: Nausícaa}

\propuesta{Puntuación de locución}{src/odisea5_nausicaa.tex}{línea 13}{La locución explicativa \emph{o sea} debe ir seguida de coma.}{También nos dice «niña en presencia y figura igual a las siemprevivas», o sea una muchachita tan hermosa como una deidad.}{También nos dice «niña en presencia y figura igual a las siemprevivas», o sea\camb{,} una muchachita tan hermosa como una deidad.}

\propuesta{Locución contaminada}{src/odisea5_nausicaa.tex}{línea 142}{\emph{En cuero picado} parece un cruce involuntario entre \emph{en cueros} y \emph{en pelota picada}.}{Sale de entre las matas, se cubre con lo que puede y se dirige a las muchachas, casi en cuero picado.}{Sale de entre las matas, se cubre con lo que puede y se dirige a las muchachas, \camb{casi en cueros}.}

\propuesta{Construcción correlativa}{src/odisea5_nausicaa.tex}{línea 222}{En la lengua cuidada, \emph{no solo} se corresponde con \emph{sino también}.}{Nausícaa no solo es valiente, también es generosa.}{Nausícaa no solo es valiente, \camb{sino también} generosa.}

\propuesta{Repetición léxica}{src/odisea5_nausicaa.tex}{línea 240}{\emph{Así} aparece dos veces en dos oraciones consecutivas.}{Pero Ulises, pudoroso, le pide a las muchachas que le dejen arreglarse solo. Así lo hace, un poco apartado, y Atenea aprovecha para maquillarlo divinalmente, así que cuando sale a la vista, está hecho un brazo de mar.}{Pero Ulises, pudoroso, le pide a las muchachas que le dejen arreglarse solo. Así lo hace, un poco apartado, y Atenea aprovecha para maquillarlo divinalmente\camb{. Cuando} sale a la vista, está hecho un brazo de mar.}

\propuesta{Puntuación}{src/odisea5_nausicaa.tex}{línea 258}{No se separa con coma la consecutiva formada por \emph{tanto que}.}{El patito (o más bien el ganso) feo se ha convertido en un apuesto galán por obra y gracia de Atenea, que de paso habrá aprovechado para quitarle alguna cana; tanto, que a la princesa le da por pensar que uno así no sería mal marido.}{El patito (o más bien el ganso) feo se ha convertido en un apuesto galán por obra y gracia de Atenea, que de paso habrá aprovechado para quitarle alguna cana; tanto\camb{ que} a la princesa le da por pensar que uno así no sería mal marido.}

\propuesta{Elipsis poco clara}{src/odisea5_nausicaa.tex}{línea 275}{Tras \emph{el tema del matrimonio sigue}, la coordinación \emph{pero también el de la maledicencia} deja el verbo excesivamente elidido.}{El tema del matrimonio con el extranjero sigue dando vueltas en el relato, pero también el de la maledicencia, y eso que todavía no se habían inventado las redes sociales.}{El tema del matrimonio con el extranjero sigue dando vueltas en el relato, \camb{y también lo hace} el de la maledicencia, y eso que todavía no se habían inventado las redes sociales.}

\propuesta{Posesivo innecesario}{src/odisea5_nausicaa.tex}{línea 305}{En español se cierran \emph{los ojos}, no normalmente \emph{tus ojos}; el posesivo parece un calco del inglés.}{Cada vez que cierras tus ojos al tenderte en la arena,}{Cada vez que cierras \camb{los} ojos al tenderte en la arena,}

\section{Sexta parte: Esqueria}

\propuesta{Correlación temporal}{src/odisea6_esqueria.tex}{línea 179}{La acomodación de Telémaco ocurrió en un capítulo anterior al momento narrado; el pluscuamperfecto fija mejor la anterioridad.}{Terminada la conversación, Arete en persona se ocupa de acomodar a Odiseo ---por supuesto bajo el soportal, el mismo lugar donde Helena ha acomodado a Telémaco--- y, una vez que ha cumplido con su deber, Alcínoo y Arete hacen exactamente lo mismo que Helena y Menelao.}{Terminada la conversación, Arete en persona se ocupa de acomodar a Odiseo ---por supuesto bajo el soportal, el mismo lugar donde Helena \camb{había acomodado} a Telémaco--- y, una vez que ha cumplido con su deber, Alcínoo y Arete hacen exactamente lo mismo que Helena y Menelao.}

\propuesta{Repetición de una idea}{src/odisea6_esqueria.tex}{línea 233}{El capítulo de Helena ya ha formulado la observación de que los aqueos son muy llorones. Aquí el recordatorio entre rayas no añade información a la escena.}{Odiseo, al oírlo, no puede contener el llanto ---ya hemos dicho que los aqueos son muy llorones---:}{Odiseo, al oírlo, no puede contener el llanto\camb{:}}

\section{La Gran Guerra}

\propuesta{Repetición léxica}{src/odisea6_laguerra.tex}{línea 3}{\emph{Todavía} aparece dos veces en la misma enumeración descriptiva.}{Imagínatelo sentado junto al fuego, enjuto, su antigua corpulencia convertida en tendón y nervio, los ojos todavía feroces a pesar del temblor en las manos y el andar encorvado, la mente todavía lúcida, quizás más lúcida que nunca.}{Imagínatelo sentado junto al fuego, enjuto, su antigua corpulencia convertida en tendón y nervio, los ojos todavía feroces a pesar del temblor en las manos y el andar encorvado, la mente \camb{aún} lúcida, quizás más lúcida que nunca.}

\propuesta{Repetición léxica}{src/odisea6_laguerra.tex}{línea 25}{\emph{El juramento} se repite en la misma oración.}{Odiseo recuerda el juramento que los príncipes aqueos habían hecho y que les obligaba a defenderse los unos a los otros, el juramento que invoca Agamenón para convocarlos:}{Odiseo recuerda el juramento que los príncipes aqueos habían hecho y que les obligaba a defenderse los unos a los otros, \camb{el que} invoca Agamenón para convocarlos:}

\propuesta{Repetición de una idea}{src/odisea6_laguerra.tex}{líneas 3 y 47}{El arranque ya presenta la mente del anciano como lúcida, incluso «más lúcida que nunca». La frase introductoria al segundo poema vuelve a afirmar lo mismo.}{Viejo y amargado, Ulises no ha perdido la lucidez ni la capacidad de reírse de sí mismo.}{Viejo y amargado, Ulises \camb{no ha perdido la capacidad de reírse de sí mismo}.}

\propuesta{Locución}{src/odisea6_laguerra.tex}{línea 67}{La expresión asentada es \emph{todo músculo}, en singular.}{todo músculos y sin ningún ingenio. Pero, si eras tan inteligente,}{\camb{todo músculo} y sin ningún ingenio. Pero, si eras tan inteligente,}

\propuesta{Orden sintáctico}{src/odisea6_laguerra.tex}{líneas 75--77}{El sujeto \emph{Aquiles} queda pospuesto de forma que parece complemento de \emph{opción}.}{Quizás se debatiese si era una mejor opción\\
Aquiles, pero la verdad es que un gran héroe\\
no es tan buena inversión como un rey poderoso.}{Quizás se debatiese \camb{si Aquiles era una opción mejor},\\
pero la verdad es que un gran héroe\\
no es tan buena inversión como un rey poderoso.}

\propuesta{Repetición léxica}{src/odisea6_laguerra.tex}{líneas 204--209}{\emph{Los demás} aparece dos veces en líneas consecutivas.}{El semidiós debía cumplir su destino\\
y quedó tan abandonado como todos los demás en las costas de Ilión,\\
pero, a diferencia de los demás, no estaba interesado en el futuro.}{El semidiós debía cumplir su destino\\
y quedó tan abandonado como todos los demás en las costas de Ilión,\\
pero, a diferencia \camb{del resto}, no estaba interesado en el futuro.}

\propuesta{Régimen verbal}{src/odisea6_laguerra.tex}{líneas 240--241}{\emph{Salirse con la suya} no rige normalmente un complemento introducido por \emph{en} con este sentido.}{Tal vez sí, pero el orgullo y la lujuria\\
siempre se han salido con la suya en hombres poderosos.}{Tal vez sí, pero el orgullo y la lujuria\\
siempre \camb{se han impuesto en los hombres poderosos}.}

\section{Séptima parte: Cíclope}

\propuesta{Puntuación explicativa}{src/odisea7_ciclope.tex}{línea 243}{\emph{Y otro detalle} introduce una explicación completa; los dos miembros no deben unirse solo con coma.}{Y otro detalle, Néstor pregunta quiénes son a sus invitados después de darles de comer, es decir, después de haberles demostrado hospitalidad.}{Y otro detalle\camb{:} Néstor pregunta quiénes son a sus invitados después de darles de comer, es decir, después de haberles demostrado hospitalidad.}

\propuesta{Repetición léxica}{src/odisea7_ciclope.tex}{líneas 284--287}{\emph{Ejemplo} aparece cinco veces en dos párrafos próximos. La propuesta conserva el razonamiento y reduce dos apariciones.}{El amigo Christopher se pasa toda su película obsesionado con el desastre que supone violar la \emph{xenía}, pero el único ejemplo sustancial que nos da ---la treta del caballo--- resulta, como hemos visto, bastante endeble. [...] Y nuestro admirado director, con la misma poca fortuna con que prescinde del ejemplo más vivo de hospitalidad en la \emph{Odisea}, ignora el mejor ejemplo en el que esta se viola de manera racional y consciente.}{El amigo Christopher se pasa toda su película obsesionado con el desastre que supone violar la \emph{xenía}, pero el único \camb{caso} sustancial que nos da ---la treta del caballo--- resulta, como hemos visto, bastante endeble. [...] Y nuestro admirado director, con la misma poca fortuna con que prescinde de \camb{la encarnación más viva} de la hospitalidad en la \emph{Odisea}, ignora el mejor ejemplo en el que esta se viola de manera racional y consciente.}

\propuesta{Nombre equivocado}{src/odisea7_ciclope.tex}{línea 287}{La contraposición es entre el Polifemo de Nolan y el de Homero. \emph{El Polifemo de Odiseo} designaría, en cambio, el relato interno del propio héroe.}{El problema del cíclope de Nolan es que es un pobre hombre, un simple, que se zampa a los aqueos como quien se merienda sardinas, sin darle, aparentemente, importancia al hecho, mientras que el Polifemo de Odiseo es artero, cruel, arrogante... Y astuto.}{El problema del cíclope de Nolan es que es un pobre hombre, un simple, que se zampa a los aqueos como quien se merienda sardinas, sin darle, aparentemente, importancia al hecho, mientras que el Polifemo de \camb{Homero} es artero, cruel, arrogante... Y astuto.}

\propuesta{Puntuación}{src/odisea7_ciclope.tex}{línea 320}{La prótasis condicional debe separarse de la oración principal.}{Ulises piensa en matarlo y se da cuenta de que si lo hace quedan atrapados en la cueva, ya que no tienen forma de mover el enorme peñasco.}{Ulises piensa en matarlo y se da cuenta de que si lo hace\camb{,} quedan atrapados en la cueva, ya que no tienen forma de mover el enorme peñasco.}

\propuesta{Construcción verbal}{src/odisea7_ciclope.tex}{línea 385}{\emph{Garantizando que lo deja} mezcla el resultado ya producido con una subordinada que aquí pide otra construcción.}{En el original, Ulises lo emborracha primero, garantizando que lo deja fuera de juego y puede maniobrar a sus anchas.}{En el original, Ulises lo emborracha primero, \camb{asegurándose de dejarlo} fuera de juego y poder maniobrar a sus anchas.}

\propuesta{Puntuación de locución}{src/odisea7_ciclope.tex}{línea 409}{La locución explicativa \emph{o sea} debe ir seguida de coma.}{o sea Nadie, Ninguno.}{o sea\camb{,} Nadie, Ninguno.}

\propuesta{Repetición léxica}{src/odisea7_ciclope.tex}{línea 519}{\emph{Pedirle}/\emph{pedirles} aparece en la misma frase y entorpece la comparación.}{Polifemo no podía pedirle a Poseidón ---que presumiblemente no estaba mirando en ese momento--- venganza contra Nadie, como no pudo pedirles a los cíclopes ayuda cuando Ninguno lo mataba con artería y no con fuerza.}{Polifemo no podía pedirle a Poseidón ---que presumiblemente no estaba mirando en ese momento--- venganza contra Nadie, como tampoco pudo \camb{solicitar ayuda a} los cíclopes cuando Ninguno lo mataba con artería y no con fuerza.}

\propuesta{Puntuación de diálogo}{src/odisea7_ciclope.tex}{líneas 526--529}{\emph{Gritan} es una acotación del narrador y debe distinguirse de las palabras de los nietos.}{¡Polifemo! Gritan. Por favor, cuéntanos, abuelo,\\
cómo derrotaste al monstruo.}{¡Polifemo! \camb{---gritan---.} Por favor, cuéntanos, abuelo,\\
cómo derrotaste al monstruo.}

\propuesta{Verbo impropio}{src/odisea7_ciclope.tex}{línea 568}{Un relato no suele \emph{argumentar} un dato; lo afirma, sostiene o alega.}{Claro, tu relato argumenta que necesitabas comida y agua,}{Claro, tu relato \camb{afirma} que necesitabas comida y agua,}

\propuesta{Puntuación de diálogo}{src/odisea7_ciclope.tex}{líneas 589--591}{Las preguntas y respuestas reproducidas necesitan marcas que permitan atribuir cada voz sin ambigüedad.}{Esa parte de la historia es cierta, al menos.\\
Los otros llamaron: ¿Quién te ha cegado? Y el pobre Polifemo gritó.\\
Nadie lo hizo.}{Esa parte de la historia es cierta, al menos.\\
Los otros llamaron: \camb{«¿Quién te ha cegado?».} Y el pobre Polifemo gritó\camb{:}\\
\camb{«Nadie lo hizo».}}

\section{Eolo: La bolsa de los vientos}

\propuesta{Redundancia}{src/odisea8_eolo.tex}{línea 138}{\emph{Anuncia} y \emph{presagio} expresan dos veces la misma relación.}{Lo que Odiseo cuenta entonces a los feacios anuncia ya un oscuro presagio.}{Lo que Odiseo cuenta entonces a los feacios \camb{constituye ya} un oscuro presagio.}

\section{Octava parte: Circe}

\propuesta{Régimen verbal}{src/odisea8_circe.tex}{línea 21}{\emph{Animar} lleva complemento directo; la reescritura evita el leísmo plural.}{Al día siguiente les anima a que vayan a explorar, pero la tripulación, después de habérselas visto con cíclopes y lestrigones, no está para muchas aventuras.}{Al día siguiente \camb{los anima a ir} a explorar, pero la tripulación, después de habérselas visto con cíclopes y lestrigones, no está para muchas aventuras.}

\propuesta{Ambigüedad sintáctica}{src/odisea8_circe.tex}{línea 95}{\emph{Le pide quedarse} puede significar que Euríloco pide a Odiseo que se quede.}{Ulises se dispone a explorar y le pide a su lugarteniente que le acompañe, pero este, muerto de miedo, le pide quedarse.}{Ulises se dispone a explorar y le pide a su lugarteniente que le acompañe, pero este, muerto de miedo, \camb{le pide que lo deje quedarse}.}

\propuesta{Repetición de una idea}{src/odisea8_circe.tex}{línea 291}{La primera oración resume lo que la cita inmediata vuelve a decir casi con las mismas palabras.}{Circe manda a su amante a hablar con Tiresias en el Hades, ese sitio horrible donde iban las almas de los mortales para convertirse en sombras sin entendimiento. Los versos son espeluznantes: «solo a este ciego agorero, / cuando murió, le mantuvo Perséfone su entendimiento / y lucidez, en un mundo de sombras en revoloteo».}{La \camb{justificación de Circe es espeluznante}: «solo a este ciego agorero, / cuando murió, le mantuvo Perséfone su entendimiento / y lucidez, en un mundo de sombras en revoloteo».}

\propuesta{Repetición y concordancia temporal}{src/odisea8_circe.tex}{líneas 472--476}{\emph{Lamentaba}/\emph{se lamentó}/\emph{se lamentó} se acumula y alterna imperfecto con perfecto.}{Y la canción también lamentaba a Eurídice,\\
que no pudo ser rescatada del Inframundo,\\
y se lamentó por la esposa de Héctor, esclavizada,\\
y se lamentó por todas esas otras mujeres que fueron tomadas,\\
por carniceros como tú.}{Y la canción también lamentaba a Eurídice,\\
que no pudo ser rescatada del Inframundo,\\
\camb{a la esposa de Héctor, esclavizada,}\\
\camb{y a todas esas otras mujeres que fueron tomadas}\\
por carniceros como tú.}

\propuesta{Puntuación de discurso directo}{src/odisea8_circe.tex}{líneas 492--496}{\emph{Su voz dijo} introduce palabras literales, que conviene marcar con dos puntos y comillas.}{y un ofrecimiento de perdón. Su voz dijo\\
monstruo o no, eres solo un hombre, y todos los hombres,

lo merezcan o no,\\
pueden ser redimidos.}{y un ofrecimiento de perdón. Su voz dijo\camb{:}\\
\camb{«Monstruo} o no, eres solo un hombre, y todos los hombres,

lo merezcan o no,\\
pueden ser redimidos\camb{».}}

\section{Novena parte: Hades}

\propuesta{Repetición léxica}{src/odisea9_hades.tex}{línea 94}{\emph{Compasión} aparece dos veces en la misma transición.}{La escena nos reconcilia con el héroe milmañas, quizás por ser la primera vez que le vemos demostrando compasión, sufriendo genuinamente por el destino aciago de un pobre inocente. Pero a Odiseo aún le queda mucho sufrimiento por delante. La siguiente sombra en aparecer es su madre. Casi inmediatamente, Homero nos muestra que, compasión o no, Ulises es más duro de pelar que el bronce que esgrime:}{La escena nos reconcilia con el héroe milmañas, quizás por ser la primera vez que le vemos demostrando compasión, sufriendo genuinamente por el destino aciago de un pobre inocente. Pero a Odiseo aún le queda mucho sufrimiento por delante. La siguiente sombra en aparecer es su madre. Casi inmediatamente, Homero nos muestra que, \camb{piadoso o no}, Ulises es más duro de pelar que el bronce que esgrime:}

\propuesta{Repetición léxica}{src/odisea9_hades.tex}{línea 131}{\emph{Le va a cobrar caro} se repite en dos frases consecutivas.}{Poseidón territodatremante le va a cobrar caro que cegara a su hijo. O para ser precisos, le va a cobrar caro que fuera tan idiota como para identificarse después de hacerlo.}{Poseidón territodatremante le va a cobrar caro que cegara a su hijo. O\camb{,} para ser precisos, \camb{castigará que} fuera tan idiota como para identificarse después de hacerlo.}

\propuesta{Repetición léxica}{src/odisea9_hades.tex}{línea 192}{\emph{Viaje} aparece tres veces en un tramo breve, dos de ellas en la misma oración.}{El viaje del héroe no termina en Ítaca, ya que está obligado a realizar un último viaje.}{El viaje del héroe no termina en Ítaca, ya que está obligado a realizar una última \camb{travesía}.}

\propuesta{Repetición de raíz}{src/odisea9_hades.tex}{línea 210}{\emph{Desgarrador} y \emph{desgarradores} se suceden de forma muy visible.}{Ulises intenta abrazarla en vano, un pasaje todavía más desgarrador que todos los desgarradores versos que ofrece el canto más triste de la \emph{Odisea} y puede que de toda la literatura, hasta que Primo Levi escribiera \emph{Si esto es un hombre}:}{Ulises intenta abrazarla en vano, un pasaje todavía más desgarrador \camb{que cuantos ofrece} el canto más triste de la \emph{Odisea} y puede que de toda la literatura, hasta que Primo Levi escribiera \emph{Si esto es un hombre}:}

\propuesta{Repetición léxica}{src/odisea9_hades.tex}{línea 310}{\emph{Le regala} aparece dos veces en la misma oración.}{Áyax le regala a Héctor un cinturón ---que luego usará Aquiles para arrastrarlo con su carro alrededor de las murallas de Ilión--- y el troyano le regala al gigante argivo su espada, que sirve para que este se suicide.}{Áyax le regala a Héctor un cinturón ---que luego usará Aquiles para arrastrarlo con su carro alrededor de las murallas de Ilión--- y el troyano \camb{entrega} al gigante argivo su espada, que sirve para que este se suicide.}

\section{Décima parte: Agamenón y Clitemnestra}

\propuesta{Verbo impropio}{src/odisea10_agamenon.tex}{línea 5}{\emph{Escabechina} es sustantivo; el verbo correspondiente es \emph{escabechar}.}{En la \emph{Odisea}, es Egisto ---no Clitemnestra--- quien lo mata, tendiéndole una emboscada en la que escabechina también a muchos de sus hombres.}{En la \emph{Odisea}, es Egisto ---no Clitemnestra--- quien lo mata, tendiéndole una emboscada en la que \camb{escabecha} también a muchos de sus hombres.}

\propuesta{Redundancia}{src/odisea10_agamenon.tex}{línea 33}{\emph{Volver a repetirse} contiene dos veces la noción de repetición.}{La fórmula, de hecho, vuelve a repetirse en la segunda incursión que el lector hace a los infiernos, siguiendo las almas de los pretendientes, en el último canto de la épica.}{La fórmula, de hecho, \camb{se repite} en la segunda incursión que el lector hace a los infiernos, siguiendo las almas de los pretendientes, en el último canto de la épica.}

\propuesta{Redundancia}{src/odisea10_agamenon.tex}{línea 75}{\emph{Ya} y \emph{anteriormente} expresan la misma anterioridad.}{El perro homérico es el animal sin vergüenza que te mira directamente a los ojos y, por tanto, la traducción más exacta del epíteto, como ya comentamos anteriormente, es «descarada».}{El perro homérico es el animal sin vergüenza que te mira directamente a los ojos y, por tanto, la traducción más exacta del epíteto, como \camb{ya comentamos}, es «descarada».}

\propuesta{Repetición léxica}{src/odisea10_agamenon.tex}{línea 109}{\emph{También} aparece tres veces en un tramo muy corto.}{También le confirma que, en asunto de mujeres, Odiseo tiene suerte. Aunque lleva un año pasándoselo estupendamente con Circe, Penélope lo espera en casa y también le espera su hijo. Agamenón, que cae poco simpático, también se lamenta de no haber podido ver al suyo y es casi el único momento en el que nos compadecemos de él.}{También le confirma que, en asunto de mujeres, Odiseo tiene suerte. Aunque lleva un año pasándoselo estupendamente con Circe, Penélope lo espera en casa y \camb{allí lo espera también} su hijo. Agamenón, que cae poco simpático, \camb{se lamenta a su vez} de no haber podido ver al suyo y es casi el único momento en el que nos compadecemos de él.}

\propuesta{Repetición y antecedente ambiguo}{src/odisea10_agamenon.tex}{línea 134}{\emph{Versión} se repite y la última aparición puede referirse tanto a Nolan como al poema propio.}{De hecho, en mi \emph{Abandonando Ítaca}, también prefiero la versión de Nolan a la de la \emph{Odisea}. Una versión en la que, por cierto, el amigo Odiseo no sale muy bien parado.}{De hecho, en mi \emph{Abandonando Ítaca} también prefiero la versión de Nolan a la de la \emph{Odisea}. \camb{En mi relato}, por cierto, el amigo Odiseo no sale muy bien parado.}

\propuesta{Régimen verbal}{src/odisea10_agamenon.tex}{líneas 148--150}{Se culpa \emph{a alguien de algo}; \emph{culpar algo por} es un calco del inglés.}{Ese idiota de Calchas,\\
culpando al capricho de los dioses por la calma chicha,\\
pidiendo un sacrificio de sangre.}{Ese idiota de Calchas,\\
\camb{culpando a los dioses de la calma chicha},\\
pidiendo un sacrificio de sangre.}

\propuesta{Repetición de raíz}{src/odisea10_agamenon.tex}{líneas 150--153}{\emph{Sacrificio}/\emph{sacrificio}/\emph{sacrificada} se concentra en cuatro versos.}{pidiendo un sacrificio de sangre.\\
No cualquier sacrificio,\\
sino que nada menos que Ifigenia, la hija de Agamenón,\\
tenía que ser sacrificada para complacer a las deidades.}{pidiendo un sacrificio de sangre.\\
\camb{No una víctima cualquiera},\\
sino que nada menos que Ifigenia, la hija de Agamenón,\\
tenía que ser sacrificada para complacer a las deidades.}

\propuesta{Repetición léxica}{src/odisea10_agamenon.tex}{líneas 163--166}{\emph{Los dioses} aparece al final de un verso y vuelve a abrir el siguiente.}{Tal vez Calchas tenía sus propias cuentas pendientes,\\
quizás los dioses estaban realmente aburridos y querían algo de diversión.\\
Salvo que los dioses ---tu vejez lo ha revelado---,\\
solo existen en la imaginación de los hombres,}{Tal vez Calchas tenía sus propias cuentas pendientes,\\
quizás \camb{los inmortales} estaban realmente aburridos y querían algo de diversión.\\
Salvo que los dioses ---tu vejez lo ha revelado---,\\
solo existen en la imaginación de los hombres,}

\propuesta{Repetición léxica}{src/odisea10_agamenon.tex}{líneas 185--188}{\emph{Caras}/\emph{cara} se repite de inmediato; además, el poema enumera luego una segunda y una tercera.}{Ahora recuerdas las caras de Ifigenia,\\
la cara que viste cuando irrumpió en la sala,\\
preguntando por su amado padre, creyendo aún}{Ahora recuerdas las caras de Ifigenia\camb{.}\\
\camb{La primera, cuando} irrumpió en la sala,\\
preguntando por su amado padre, creyendo aún}

\section{Undécima parte: Sirenas}

\propuesta{Concordancia temporal}{src/odisea11_sirenas.tex}{líneas 184--186}{La serie arranca en condicional y pasa sin motivo al imperfecto.}{Podría ser por la muerte de Héctor, o la muerte de Aquiles,\\
podía recordar las naves perdidas o invocar las llamas que incendiaban Troya,\\
lamentar el destino de Patroclo o añorar la belleza de Helena,}{\camb{Podía} ser por la muerte de Héctor, o la muerte de Aquiles,\\
podía recordar las naves perdidas o invocar las llamas que incendiaban Troya,\\
lamentar el destino de Patroclo o añorar la belleza de Helena,}

\propuesta{Construcción enfática}{src/odisea11_sirenas.tex}{línea 191}{La construcción \emph{es de... de lo que} duplica la preposición.}{Y es, justamente, de ese sentido de lo sagrado de lo que se alimentan las sirenas:}{Y es justamente de ese sentido de lo sagrado \camb{del que} se alimentan las sirenas:}

\propuesta{Puntuación}{src/odisea11_sirenas.tex}{líneas 194--197}{La segunda oración explica la primera y no debe quedar unida a ella solo por coma.}{Quizás esto es lo que hizo que las sirenas\\
fuesen tan temibles y tan peligrosas,\\
cantaban simplemente canciones de los hombres\\
con acordes más puros.}{Quizás esto es lo que hizo que las sirenas\\
fuesen tan temibles y tan peligrosas\camb{:}\\
cantaban simplemente canciones de los hombres\\
con acordes más puros.}

\section{Duodécima parte: La matanza de los pretendientes}

\propuesta{Régimen de participio}{src/odisea12_pretendientes.tex}{línea 5}{Con valor de orden temporal, la forma asentada es \emph{seguido de}; \emph{seguido por} introduce normalmente al agente que sigue.}{Un solo vocablo para nombrar el asesinato a sangre fría de ciento ocho jóvenes, seguido por una implacable ejecución.}{Un solo vocablo para nombrar el asesinato a sangre fría de ciento ocho jóvenes, seguido \camb{de} una implacable ejecución.}

\propuesta{Repetición de una idea}{src/odisea12_pretendientes.tex}{líneas 63 y 77}{La condición de adivino de Leodes se explica en la línea 63 y vuelve a presentarse catorce líneas después como algo ya dicho. Puede enlazarse sin repetir el acto de explicación.}{Ya hemos dicho que Leodes tiene dotes de adivino, así que más les habría valido a los pretendientes hacerle caso cuando dice:}{\camb{Como Leodes tiene dotes de adivino}, más les habría valido a los pretendientes hacerle caso cuando dice:}

\propuesta{Repetición léxica}{src/odisea12_pretendientes.tex}{línea 122}{\emph{Detalle} aparece dos veces en el mismo párrafo.}{El detalle no tiene desperdicio. [...] Cuando todos están bastante borrachos ---un detalle importante, ya que es más fácil abatir al enemigo si está ebrio, sin olvidar que el héroe ya ha aplicado ese truco con el cíclope---, Odiseo pone en marcha su treta:}{El detalle no tiene desperdicio. [...] Cuando todos están bastante borrachos ---una \camb{circunstancia} importante, ya que es más fácil abatir al enemigo si está ebrio, sin olvidar que el héroe ya ha aplicado ese truco con el cíclope---, Odiseo pone en marcha su treta:}

\propuesta{Repetición léxica}{src/odisea12_pretendientes.tex}{línea 188}{\emph{Importante}/\emph{importa} y \emph{también} se acumulan en la explicación.}{Y esto es importante porque su intervención favorece los planes de Ulises sin ella saberlo y también lo es porque muestra su generosidad y, sobre todo, \emph{su agencia}.}{\camb{Esto importa} porque su intervención favorece los planes de Ulises sin ella saberlo. \camb{Además, muestra} su generosidad y, sobre todo, \emph{su agencia}.}

\propuesta{Construcción recargada}{src/odisea12_pretendientes.tex}{línea 418}{La frase acumula \emph{una escena}, \emph{en la secuencia} y \emph{en Nolan}.}{Hay una escena en la secuencia de la matanza, en Nolan, que da casi risa.}{Hay una escena \camb{de la matanza de Nolan} que da casi risa.}

\propuesta{Régimen verbal}{src/odisea12_pretendientes.tex}{línea 446}{\emph{Mandar a alguien para que haga} resulta menos preciso que la construcción directa con \emph{que}.}{Cuando Melantio se dirige al almacén para seguir trayendo armas, Ulises manda a Eumeo y Filetio para que lo intercepten.}{Cuando Melantio se dirige al almacén para seguir trayendo armas, Ulises manda a Eumeo y Filetio \camb{que lo intercepten}.}

\propuesta{Verbo impropio}{src/odisea12_pretendientes.tex}{línea 534}{Una interpretación no \emph{circula en} una obra; aquí la idea la recorre o está presente en ella.}{Es una interpretación que circula, como hemos visto, en toda la \emph{Odisea} y, en general, en toda la épica griega.}{Es una interpretación que \camb{recorre}, como hemos visto, toda la \emph{Odisea} y, en general, toda la épica griega.}

\propuesta{Repetición léxica}{src/odisea12_pretendientes.tex}{línea 544}{\emph{Quizás} aparece tres veces en tres oraciones consecutivas. La incertidumbre se conserva con una sola alternancia.}{A la inversa, quizás los humanos también «echan la culpa» a los olímpicos cuando la suerte les favorece. Así que quizás no hace falta recurrir a la diosa ojiglauca para frenar las lanzas. Quizás, simplemente, Odiseo es un tipo afortunado, como se dirá a sí mismo, años más tarde, viejo y desengañado en Ítaca:}{A la inversa, quizás los humanos también «echan la culpa» a los olímpicos cuando la suerte les favorece. \camb{Por tanto, puede que no haga falta} recurrir a la diosa ojiglauca para frenar las lanzas. \camb{Tal vez}, simplemente, Odiseo es un tipo afortunado, como se dirá a sí mismo, años más tarde, viejo y desengañado en Ítaca:}

\section{Decimotercera parte: El asesinato de las niñas}

\propuesta{Relativo temporal}{src/odisea13_ninas.tex}{línea 3}{En registro cuidado, un antecedente temporal explícito pide \emph{el día en que}.}{Recuerdo como si fuera ayer el día que terminé el canto XXII.}{Recuerdo como si fuera ayer el día \camb{en que} terminé el canto XXII.}

\propuesta{Repetición léxica}{src/odisea13_ninas.tex}{línea 65}{\emph{Es decir} aparece dos veces en tres frases consecutivas.}{Es decir, nada de matarlas a espada como han hecho con los pretendientes. Después de todo, estos eran nobles y varones, es decir, se merecían un pasaje en primera al Hades.}{Es decir, nada de matarlas a espada como han hecho con los pretendientes. Después de todo, estos eran nobles y varones\camb{:} se merecían un pasaje en primera al Hades.}

\propuesta{Falta de antecedente}{src/odisea13_ninas.tex}{líneas 140--145}{El pasaje comentado que presentaba a Pentesilea ya no forma parte del capítulo. Por eso \emph{la reina de las amazonas} aparece aquí sin nombre ni presentación.}{Pero no todos se engañan:

En cambio, ni Héctor ni la reina de las amazonas\\
esperaban derrotar al guerrero invulnerable}{Pero no todos se engañan\camb{. Pentesilea, reina de las amazonas, tampoco:}

En cambio, ni Héctor ni la reina de las amazonas\\
esperaban derrotar al guerrero invulnerable}

\propuesta{Posesivo innecesario}{src/odisea13_ninas.tex}{línea 170}{En español se sacude \emph{la cabeza}, no normalmente \emph{sus cabezas}; el posesivo parece un calco del inglés.}{Amigos y sirvientes por igual\\
sacudirían sus cabezas y pensarían en un viejo que ha perdido el sentido.}{Amigos y sirvientes por igual\\
sacudirían \camb{la cabeza} y pensarían en un viejo que ha perdido el sentido.}

\section{Decimocuarta parte: Penélope}

\propuesta{Construcción disyuntiva}{src/odisea14_penelope.tex}{línea 112}{Los dos términos de la interrogación indirecta no comparten la misma estructura: \emph{si es... el que ha vuelto o se trata...}.}{Penélope sigue en sus trece y dispone de recursos para averiguar, en privado, si es de verdad su marido el que ha vuelto o se trata de un divinal impostor.}{Penélope sigue en sus trece y dispone de recursos para averiguar, en privado, \camb{si quien ha vuelto es de verdad su marido o un divinal impostor}.}

\propuesta{Construcción incompleta}{src/odisea14_penelope.tex}{línea 286}{\emph{O quizás más difícil} necesita un núcleo comparativo antes de introducir la segunda obligación.}{Ulises le hará a su esposa el relato de sus aventuras y ella tendrá que fingir que no le importan el año con Circe ni los siete años con Calipso. O quizás más difícil, tendrá que aprender a convivir con un extraño.}{Ulises le hará a su esposa el relato de sus aventuras y ella tendrá que fingir que no le importan el año con Circe ni los siete años con Calipso. O, \camb{lo que quizás sea más difícil}, tendrá que aprender a convivir con un extraño.}

\section{Decimoquinta parte: Homero}

\propuesta{Formación impropia con prefijo}{src/odisea15_homero.tex}{línea 5}{\emph{Ex-alegre} no es una formación normativa: \emph{ex-} se une a sustantivos o adjetivos sustantivados que designan una condición anterior.}{No es que Homero se interese mucho por la ex-alegre pandilla de gorrones.}{No es que Homero se interese mucho por la \camb{antaño alegre} pandilla de gorrones.}

\propuesta{Puntuación del poema}{src/odisea15_homero.tex}{líneas 20--23}{Las comas separan \emph{tanto} de su consecutiva, \emph{costumbre} de su complemento y \emph{antes} del complemento que lo precisa.}{A papá le preocupaba que aprendiera a conducir,\\
le preocupaba tanto, que tomamos la costumbre,\\
de hacer prácticas, mucho antes,\\
de tener la edad para sacarme el carnet.}{A papá le preocupaba que aprendiera a conducir,\\
le preocupaba tanto\camb{ que} tomamos la costumbre\\
de hacer prácticas mucho antes\\
de tener la edad para sacarme el carnet.}

\propuesta{Puntuación del poema}{src/odisea15_homero.tex}{líneas 35--41}{Una coma separa sujeto y verbo en \emph{lo suyo era}; otra deja aislado el complemento de \emph{tenía}. La unión de dos oraciones completas con coma puede aclararse con punto.}{Parábamos en pequeños cafés,\\
donde charlábamos largamente, después de todo,\\
éramos socios al volante y lo suyo,\\
era intercambiar historias.

Y mi padre tenía,\\
tantas historias que contarme.}{Parábamos en pequeños cafés,\\
donde charlábamos largamente\camb{.} Después de todo,\\
éramos socios al volante y lo suyo\\
era intercambiar historias.

Y mi padre tenía\\
tantas historias que contarme.}

\propuesta{Puntuación del poema}{src/odisea15_homero.tex}{líneas 43--54}{Varias comas separan el verbo de sus complementos.}{Cuando saqué el carnet siguió insistiendo,\\
en que practicar era imprescindible,\\
declarándose todavía francamente inquieto,\\
con mi obvia falta de reflejos.

Era una extraña obsesión, de hecho,\\
yo había sacado el carnet a la primera,\\
y después de tanta práctica, sabía muy bien\\
que me sobraba habilidad. No importaba,

papá declaró que todavía eran necesarias,\\
veinte mil leguas de viaje.}{Cuando saqué el carnet siguió insistiendo\\
en que practicar era imprescindible,\\
declarándose todavía francamente inquieto\\
con mi obvia falta de reflejos.

Era una extraña obsesión\camb{.} De hecho,\\
yo había sacado el carnet a la primera\\
y\camb{,} después de tanta práctica, sabía muy bien\\
que me sobraba habilidad. No importaba\camb{:}

papá declaró que todavía eran necesarias\\
veinte mil leguas de viaje.}

\propuesta{Puntuación del poema}{src/odisea15_homero.tex}{líneas 56--77}{Se acumulan comas entre verbo y complemento, sujeto y verbo, y preposición y término. En \emph{Ayer, pasé} la coma separa un adverbio breve del verbo.}{Así que las hicimos. Viajábamos a veces,\\
al pueblo veraniego, pernoctábamos,\\
en la casita que teníamos allí,\\
antes de volver a casa al día siguiente.

Papá cancelaba sus clases,\\
anulaba compromisos, se saltaba citas,\\
todo fuera por enseñarme,\\
a conducir, decía, como Dios manda.

Y mamá, de repente,\\
parecía igualmente preocupada,\\
y declaraba que aquellas excursiones,\\
eran absolutamente necesarias.

Ayer, pasé el día junto a mi hijo,\\
en otro largo viaje, conduce tan bien\\
como solía hacerlo yo, y aun así,\\
prolongaría ese viaje hasta Finisterre.

Viajaría con él veinte mil leguas, hasta llegar,\\
al final del Arco Iris, donde papá nos espera.}{Así que las hicimos. Viajábamos a veces\\
al pueblo veraniego\camb{ y} pernoctábamos\\
en la casita que teníamos allí\\
antes de volver a casa al día siguiente.

Papá cancelaba sus clases,\\
anulaba compromisos, se saltaba citas,\\
todo fuera por enseñarme\\
a conducir, decía, como Dios manda.

Y mamá, de repente,\\
parecía igualmente preocupada\\
y declaraba que aquellas excursiones\\
eran absolutamente necesarias.

Ayer pasé el día junto a mi hijo\\
en otro largo viaje\camb{.} Conduce tan bien\\
como solía hacerlo yo, y aun así,\\
prolongaría ese viaje hasta Finisterre.

Viajaría con él veinte mil leguas, hasta llegar\\
al final del Arco Iris, donde papá nos espera.}

\propuesta{Repetición léxica}{src/odisea15_homero.tex}{línea 82}{\emph{Darte} aparece dos veces dentro del mismo inciso.}{Quizás sea un prejuicio mío ---el sesgo emocional que produce darte cuenta de que tu padre ha cancelado una reunión importante para darte una clase de conducir que no necesitas--- lo que me deja frío frente a unos versos formalmente magníficos pero tan formales que uno no acaba de ver el sentimiento:}{Quizás sea un prejuicio mío ---el sesgo emocional que produce \camb{descubrir} que tu padre ha cancelado una reunión importante para darte una clase de conducir que no necesitas--- lo que me deja frío frente a unos versos formalmente magníficos pero tan formales que uno no acaba de ver el sentimiento:}

\propuesta{Redundancia}{src/odisea15_homero.tex}{línea 195}{\emph{Tiene cien años} ya expresa la antigüedad.}{La noción de que ambos poemas épicos se derivan de una amplia ---y muy antigua--- tradición oral tiene ya cien años de antigüedad.}{La noción de que ambos poemas épicos se derivan de una amplia ---y muy antigua--- tradición oral \camb{tiene ya cien años}.}

\propuesta{Repetición léxica}{src/odisea15_homero.tex}{línea 205}{\emph{Fórmulas} aparece cuatro veces en un párrafo. La propuesta reserva el término para la definición de Parry y para la serie final.}{Todo un armazón de trucos nemotécnicos y fórmulas prefabricadas disponible para que el bardo pueda memorizar e improvisar. De hecho, Parry llamó a estos patrones fijos, precisamente, «fórmulas». Ya hemos visto algunas, ahí van algunas más:}{Todo un armazón de trucos nemotécnicos y \camb{expresiones} prefabricadas disponible para que el bardo pueda memorizar e improvisar. De hecho, Parry llamó a estos patrones fijos, precisamente, «fórmulas». Ya hemos visto algunas\camb{; ahí van otras}:}

\propuesta{Repetición léxica}{src/odisea15_homero.tex}{línea 233}{\emph{Varios} se repite en la misma frase.}{El aparente batiburrillo de estilos y capas dialectales, que sugería que el poema había sido compuesto por varios bardos desparramados en varias épocas.}{El aparente batiburrillo de estilos y capas dialectales, que sugería que el poema había sido compuesto por varios bardos desparramados en \camb{distintas} épocas.}

\section{Vigésima parte: Marcharse de Ítaca}

\propuesta{Repetición léxica}{src/odisea_marcharse_de_itaca.tex}{línea 5}{\emph{Acabábamos}/\emph{acababa} abre dos frases consecutivas.}{Acabábamos de cumplir los dieciséis años. Acababa de marcharme de Alba, donde había llegado a los diez, justo cuando parecía que el pueblo empezaba a aceptarme.}{Acabábamos de cumplir los dieciséis años. \camb{Me había marchado} de Alba, donde había llegado a los diez, justo cuando parecía que el pueblo empezaba a aceptarme.}

\propuesta{Topónimo}{src/odisea_marcharse_de_itaca.tex}{línea 7}{La forma oficial y ortográficamente completa del topónimo es \emph{Sant Feliu de Guíxols}.}{Los turistas embarcaban en las golondrinas que cubrían el trayecto entre Blanes y San Feliú de Guixols, deteniéndose un par de horas en Tossa de Mar.}{Los turistas embarcaban en las golondrinas que cubrían el trayecto entre Blanes y \camb{Sant Feliu de Guíxols}, deteniéndose un par de horas en Tossa de Mar.}

\propuesta{Repetición léxica}{src/odisea_marcharse_de_itaca.tex}{línea 28}{\emph{Encontré} aparece tres veces en la evocación de la Ítaca industrial. Basta variar una aparición para descargar la serie sin romper su paralelismo.}{Pero quizás aquella Ítaca industrial de mi adolescencia no era sino el país de los feacios, donde encontré la hospitalidad que no conocía aún, donde conocí no a una, sino a muchas Nausícaas, donde, más adelante, encontré y luego abandoné a Penélope, rompiendo su corazón y el mío.}{Pero quizás aquella Ítaca industrial de mi adolescencia no era sino el país de los feacios, donde \camb{hallé} la hospitalidad que no conocía aún, donde conocí no a una, sino a muchas Nausícaas, donde, más adelante, encontré y luego abandoné a Penélope, rompiendo su corazón y el mío.}

\propuesta{Redundancia adversativa}{src/odisea_marcharse_de_itaca.tex}{líneas 122--123}{\emph{Irónicamente} y \emph{sin embargo} cumplen aquí la misma función adversativa. La propuesta mantiene los dos versos, pero da contenido propio al primero.}{Irónicamente,\\
sin embargo, ella te amaba,}{\camb{Contra todo pronóstico},\\
ella te amaba,}

\propuesta{Concordancia temporal}{src/odisea_marcharse_de_itaca.tex}{líneas 128--129}{La mirada y la visión están narradas en imperfecto; \emph{merece} introduce un presente que no parece deliberado.}{Veías un monstruo\\
que no merece ser amado.}{Veías un monstruo\\
que no \camb{merecía} ser amado.}

\propuesta{Tiempo verbal}{src/odisea_marcharse_de_itaca.tex}{línea 143}{Con el adverbio \emph{ahora}, el perfecto compuesto o el presente resultan naturales; el perfecto simple sitúa el hecho en un pasado cerrado.}{Ahora llegó el momento\\
para que intentes encontrar tu camino}{Ahora \camb{ha llegado} el momento\\
para que intentes encontrar tu camino}

\propuesta{Construcción incompleta}{src/odisea_marcharse_de_itaca.tex}{líneas 168--172}{\emph{Has llegado a pensar} necesita aquí la conjunción \emph{que} para introducir el contenido pensado.}{Tal vez todavía ella tenga la droga\\
que una vez te dejó caminar por el Hades\\
---quizás, has llegado a pensar,\\
el Hades no es más que las profundidades\\
de tu mente atormentada---.}{Tal vez todavía ella tenga la droga\\
que una vez te dejó caminar por el Hades\\
---quizás has llegado a pensar \camb{que}\\
el Hades no es más que las profundidades\\
de tu mente atormentada---.}

\clearpage
\section{Recapitulación}

El informe contiene \textbf{\thecambio{} cambios propuestos}. Ninguno se ha aplicado al corpus. La mayor concentración de incidencias normativas aparece en la puntuación de «Lecciones de conducir»; se ha agrupado por bloques para que la regularización pueda aprobarse sin convertir cada coma en un cambio independiente.

No se han marcado como problemas, entre otras recurrencias conscientes, la triple formulación de «El problema» en la introducción, las cadenas de \emph{quizás} que expresan conjetura cuando conservan un ritmo claro, el retorno de «No es de extrañar que el cine tema a Homero», los estribillos del capítulo XX, las fórmulas homéricas ni las repeticiones internas de las traducciones citadas. Tampoco se propone suprimir verso alguno.

\end{document}
